\section{Related Work}\label{sec:related_work}

\begin{figure*}[t]
\centering
\includegraphics[width=0.95\textwidth]{figures/fig_pipeline.pdf}
\caption{ContamCal three-stage pipeline. Stage~1 fuses training-free detection signals into per-item exposure scores via Reciprocal Rank Fusion. Stage~2 embeds these scores in a contamination-augmented 2PL IRT model, recovering clean ability estimates with credible intervals. Stage~3 classifies the resulting dose-response trajectories into canonical patterns.}
\label{fig:pipeline}
\end{figure*}

\paragraph{Contamination detection.}
Detection methods flag contaminated items via token-level surprisal \citep{shi_2023_minkprob,zhang_2024_minkpp}, accuracy gaps on paraphrases or rephrased items \citep{dekoninck2024constat,yang2023rethinking}, kernel divergence \citep{choi_2025_kds}, completion probes \citep{dong_2024_cdd,golchin2024time}, and parallel test sets \citep{zhang_2024_gsm1k}; see \citet{chen2025benchmarking,singh2024evaluation,sainz2023nlp} for surveys. \citet{balloccu2024leak} document contamination malpractices in closed-source LLMs; \citet{bordt2025forget} argue that contamination effects are often negligible in practice. Earlier work characterized contamination as memorization \citep{carlini2023quantifying,magar_2022_contamination,li_2024_task_contamination,xu_2024_leakage}; \citet{tao2025detecting} and \citet{kocyigit2025impact} address RL post-training. These methods perform binary classification without quantifying effect sizes.

\paragraph{Contamination quantification and mitigation.}
On mitigation, \citet{sun2025emperor} find none of twenty approaches reliably restores clean performance; DCR \citep{xu2025dcr} estimates decontaminated scores but without uncertainty quantification. ContamCal consumes detection signals as inputs and produces calibrated ability with credible intervals. Three concurrent lines of work directly address contamination quantification. \citet{han2026quantifying} apply observational differential item functioning (DIF) to detect contamination without controlled dosage, achieving detection but not dose-response profiling. \citet{schaeffer2026causally} perform controlled pretraining-stage injection with larger seed counts ($n\!=\!10$), establishing causal contamination effects at the pretraining level; our $n\!=\!3$ seed design trades statistical power for broader family coverage. Our work complements both: like Schaeffer et al., we use controlled injection for causal identification, but at the SFT stage with precise dosage control; like Han et al., we employ IRT-based modeling, but with ground-truth contamination labels enabling dose-response characterization.

\paragraph{IRT for LLM evaluation.}
Item Response Theory \citep{lord_1980_irt,hambleton_1991_irt} has gained traction for LLM evaluation \citep{liang_2023_helm}: \citet{rodriguez2021evaluation}, TinyBenchmarks \citep{polo_2024_tinybenchmarks}, Growing Pains \citep{habba_2025_growing_pains}, HiBayES \citep{luettgau_2025_hibayes}, PSN-IRT \citep{zhou_2025_psnirt}; see \citet{lalor_2024_irt_nlp} for a tutorial and \citet{gu2024olmes} for evaluation standardization. Most directly related, \citet{han2026quantifying} apply group-based DIF \citep{holland_1993_dif} to quantify item-level contamination effects. ContamCal differs by using continuous exposure covariates (explanatory DIF; \citealp{deboeck_2004_explanatory_irt}), providing calibrated ability with credible intervals, and domain-level diagnosis ($\gamma_d$); the approaches are complementary (item-level DIF for benchmark curation, model-level calibration for score interpretation). Our $\gamma$ parameter is formally an instance of explanatory DIF \citep{deboeck_2004_explanatory_irt,eckerly_2017_item_compromise}: the contamination covariate differentially affects response probability, analogous to item compromise models in educational measurement \citep{sinharay_2017_preknowledge,holland_1993_dif,camilli_1994_dif,swaminathan_1990_dif}. We use ``calibration'' in the psychometric sense of adjusting scores to approximate uncontaminated performance, distinct from probability calibration \citep{guo_2017_calibration}.
